% !PRINTED 454
\weberpage{454}

Dieser Satz, der uns in der Folge noch sehr wichtige Dienste leisten wird, ist fast selbstverständlich; denn der grösste gemeinschaftliche Theiler zweier Functionen $F(x)$ und $f(x)$ lässt sich durch rationale Rechnung finden, und ist daher auch in $\Omega$ enthalten. Da nun aber $f(x)$ keinen in $\Omega$ enthaltenen Theiler hat als sich selbst oder eine Constante (d.\,h. eine in $\Omega$ enthaltene Grösse), so muss also dieser grösste gemeinschaftliche Theiler entweder eine Constante oder $f(x)$ selbst sein.

Es folgt aus diesem Satz, dass eine irreducible Function niemals mehrfache Factoren haben kann; denn sie müsste sonst mit ihrer Ableitung $f'(x)$ einen gemeinsamen Theiler haben; also müsste $f'(x)$ durch $f(x)$ theilbar sein, was unmöglich ist, da der Grad von $f'(x)$ niedriger ist als der von $f(x)$.

Wir können dem Satze I auch den folgenden Ausdruck geben:

\gesperrt{II. Ist $f(x)$ irreducibel, und verschwindet $F(x)$ für eine Wurzel von $f(x) = 0$, so verschwindet es auch für alle anderen Wurzeln von $f(x) = 0$.}

Insbesondere können wir daraus schliessen:

\gesperrt{III. Ist $F(x)$ von niedrigerem Grade als die irreducible Function $f(x)$, und verschwindet $F(x)$ für eine Wurzel von $f(x)$, so muss $F(x)$ identisch verschwinden, d.\,h. alle Co\"efficienten von $F(x)$ müssen Null sein.}

Auch ganze rationale Functionen von mehreren Variablen, $f(x, y, z \ldots)$, deren Co\"efficienten dem Körper $\Omega$ angehören, heissen Functionen in $\Omega$. Auch diese werden als \gesperrt{reducible} und \gesperrt{irreducible} Functionen bezeichnet, je nachdem sie zerlegbar oder nicht zerlegbar sind in mehrere Factoren, die selbst Functionen in $\Omega$ sind, und wenigstens einige der Variablen $x, y, z \ldots$ wirklich enthalten.

Hier ist aber auf einen Unterschied aufmerksam zu machen.

Unter den Functionen von mehr als einer unabhängigen Veränderlichen giebt es solche, die überhaupt nicht in Factoren niedrigeren Grades, die rational von den Veränderlichen abhängen, zerlegbar sind, und wieder andere, die zwar zerlegbar sind, aber nur in solche Factoren, deren Co\"efficienten nicht in $\Omega$ enthalten sind, und endlich drittens Functionen, die in mehrere Functionen in $\Omega$ zerlegbar sind; nur die letztere Art werden wir reducibel in $\Omega$ nennen, während die der zweiten Art als \gesperrt{zerleg-}
%CONT>>
